\documentclass[10pt]{scrartcl}
\usepackage[mdthm,colorsec,noauthor,noasy]{evan}
\usepackage{handout}
\newcommand{\aug}[2]{\left[\begin{array}{@{}#1@{}}#2\end{array}\right]}
\newcommand{\vv}[1]{\mathbf{#1}}
\title{Math 54 --- Week 3}
\subtitle{Solutions}
\author{}\date{}\ihead{}
\begin{document}
\maketitle

\begin{prob}\textbf{P1.}\par\idea{The columns of $A$ are linearly independent if and only if $A$ has a pivot in every column.}
Applying $R_2\leftarrow R_2-3R_1$ and $R_3\leftarrow R_3+2R_1$, followed by $R_3\leftarrow R_3-R_2$, yields the echelon form
\[
  \begin{bmatrix} 1 & -3 & 2 \\ 0 & 1 & -1 \\ 0 & 0 & 2 \end{bmatrix},
\]
which has $3$ pivots (one in every column). Thus the columns are \textbf{linearly independent}.
\end{prob}

\begin{prob}\textbf{P2.}\par\idea{Row reduce with $h$ in place; linear dependence holds iff the third column has no pivot.}
\[
  \begin{bmatrix}1&3&-1\\-1&-5&5\\4&7&h\end{bmatrix}\xrightarrow[R_3-4R_1]{R_2+R_1}\begin{bmatrix}1&3&-1\\0&-2&4\\0&-5&h+4\end{bmatrix}\xrightarrow{R_3-\frac52R_2}\begin{bmatrix}1&3&-1\\0&-2&4\\0&0&h-6\end{bmatrix}.
\]
The vectors are linearly dependent if and only if $h-6=0$, i.e.\ $h=6$. When $h=6$, $x_3$ is free; setting $x_3=1$ gives $-2x_2+4=0\implies x_2=2$ and $x_1=-3(2)+1=-5$, yielding the linear dependence relation
\[
  -5(1,-1,4) + 2(3,-5,7) + 1(-1,5,6) = (0,0,0).
\]
\end{prob}

\begin{prob}\textbf{P3.}
(a) \textbf{Linearly dependent} by \xlink{math54}{2026-09-07-independence-linear-transformations}{thm-lay-1-8}{Thm.~1.8} ($3$ vectors in $\RR^2$, so $p=3>2=n$).
(b) \textbf{Linearly dependent} because the first vector is $-2$ times the second vector.
(c) \textbf{Linearly dependent} by \xlink{math54}{2026-09-07-independence-linear-transformations}{thm-lay-1-9}{Thm.~1.9} because the set contains the zero vector $(0,0,0)$.
(d) \textbf{Linearly independent} because neither of the two nonzero vectors is a scalar multiple of the other.
\end{prob}

\begin{prob}\textbf{P4.}\par\idea{Track $\vv e_1$ and $\vv e_2$ through both geometric transformations (\xlink{math54}{2026-09-07-independence-linear-transformations}{thm-lay-1-10}{Thm.~1.10}).}
For $\vv e_1=(1,0)$: reflection across the $x_1$-axis gives $(1,0)$, and rotation by $90^\circ$ gives $T(\vv e_1)=(0,1)$. For $\vv e_2=(0,1)$: reflection across the $x_1$-axis gives $(0,-1)$, and rotation by $90^\circ$ gives $T(\vv e_2)=(1,0)$. Thus the standard matrix is
\[
  A = \begin{bmatrix} T(\vv e_1) & T(\vv e_2) \end{bmatrix} = \begin{bmatrix} 0 & 1 \\ 1 & 0 \end{bmatrix},
\]
which acts by $T(x_1,x_2)=(x_2,x_1)$, the reflection across the line $x_2=x_1$.
\end{prob}

\begin{prob}\textbf{P5.}
\[
  \aug{rrr|r}{1&-2&3&-1\\0&1&-1&2\\2&-3&5&0}
  \xrightarrow{R_3-2R_1}
  \aug{rrr|r}{1&-2&3&-1\\0&1&-1&2\\0&1&-1&2}
  \xrightarrow[R_1+2R_2]{R_3-R_2}
  \aug{rrr|r}{1&0&1&3\\0&1&-1&2\\0&0&0&0}.
\]
The solution set of $T(\vv x)=(-1,2,0)$ is $\vv x=(3,2,0)+x_3(-1,1,1)$ ($x_3$ free). By \xlink{math54}{2026-09-07-independence-linear-transformations}{thm-lay-1-12}{Thm.~1.12}, since $A$ has no pivot in column~3, $T$ is \textbf{not one-to-one}; and since $A$ has no pivot in row~3, $T$ is \textbf{not onto} $\RR^3$.
\end{prob}

\begin{prob}\textbf{P6.}
$T(0,0)=(0,-1)\ne(0,0)$, so $T$ is \textbf{not linear}. By contrast,
\[
  S(x_1,x_2) = \begin{bmatrix} 0 & 1 \\ 3 & -1 \\ 0 & 0 \end{bmatrix}\begin{bmatrix} x_1 \\ x_2 \end{bmatrix}
\]
is a matrix transformation, so $S$ \textbf{is linear}.
\end{prob}

\begin{prob}\textbf{P7.}
Row reducing the $3\times 4$ standard matrix of $T$ to reduced echelon form yields
\[
  \begin{bmatrix} 1 & 0 & 0 & 3 \\ 0 & 1 & 0 & 0 \\ 0 & 0 & 1 & 1 \end{bmatrix}.
\]
By \xlink{math54}{2026-09-07-independence-linear-transformations}{thm-lay-1-12}{Thm.~1.12}, $T$ has a pivot in every row ($3$ pivots in $3$ rows), so $T$ is \textbf{onto} $\RR^3$; column~4 has no pivot, so $T$ is \textbf{not one-to-one} (as guaranteed by $n=4>3=m$).
\end{prob}

\begin{prob}\textbf{P8.}
By linearity,
\[
  T(3,-2)=T(3\vv e_1-2\vv e_2)=3T(\vv e_1)-2T(\vv e_2)=3(2,-1,3)-2(0,4,1)=(6,-11,7).
\] By \xlink{math54}{2026-09-07-independence-linear-transformations}{thm-lay-1-10}{Thm.~1.10}, the standard matrix of $T$ is
\[
  A = \begin{bmatrix} T(\vv e_1) & T(\vv e_2) \end{bmatrix} = \begin{bmatrix} 2 & 0 \\ -1 & 4 \\ 3 & 1 \end{bmatrix}.
\]
\end{prob}

\end{document}
